\documentclass[12pt]{article}
\usepackage[utf8]{inputenc}
\usepackage[T1]{fontenc}
\usepackage{hyperref}
\title{Wahau Kayan Language Resources for K-12 Curriculum Development}
\author{Compiled by Claude Sonnet 5.5 (claude-sonnet-5-5) for Andr\'as Kornai}
\date{}
\begin{document}
\maketitle

\section{Introduction}
Wahau Kayan (ISO 639-3 whu; Glottocode waha1239) is an endangered Austronesian language of East Kalimantan, Indonesia, spoken north of the town of Muara Wahau. The prompt's family field reads ``Bookkeeping'', which appears to be a data error. The reference description and the sources found both say Austronesian. Speaker estimates in the sources are small and disagree: Joshua Project gives about 900, and a linguistic survey gives about 500. It is distinct from Wahau Kenyah (whk). Little published material was found, and no primer, textbook, dictionary or grammar turned up.

\section{Primary Educational Resources}
None found. The language is not known to be taught in schools.

\section{Linguistic and Reference Materials}
A conference handout on language documentation in Kalimantan lists Wahau Kayan among the endangered languages of East Kalimantan: \url{https://lingdy.aa-ken.jp/wp-content/uploads/2012/01/120217_soriente_inagaki_h2.pdf}.
No dictionary, grammar or orthography description was found.

\section{Religious and Community Resources}
Joshua Project has a profile of the Wahau Kayan people, marking the language as unwritten and in need of Bible translation: \url{https://joshuaproject.net/people_groups/12593}. Its language page is \url{https://joshuaproject.net/languages/whu}.
No community organisation was identified.

\section{Digital and Audio Resources}
No audio, video or corpus resources were found.

\section{Curriculum Development Context}
The language is reported as unwritten, and almost nothing is documented. Any curriculum would first need community-led documentation, an agreed writing system, and recorded texts. The sources disagree about whether young people still speak it, so the real vitality should be checked with the community.
\end{document}